\begin{minipage}{\linewidth}
This call to the \texttt{Dwm::What::Info} constructor:
\begin{verbatim}
   Dwm::What::Info(DWM_WHAT_TYPE_EXE DWM_WHAT_TYPE_HDR, DWM_WHAT_STATUS_REL,
                   "DwmWhat", "1.0.0", "Daniel McRobb 2025", "mcplex.net");
\end{verbatim}

Would produce a 67 byte string literal at compile time with the layout shown
in \autoref{figure:InfoEx01}.  Note this figure shows what would be returned
by accessors inherited from the \texttt{Dwm::What::SegmentedLiteral} base
class (in black) and the accessors from the \texttt{Dwm::What::Info} class
(in \textcolor{blue}{blue}).

\begin{center}
  \begin{figure}[H]
\noindent\begin{tikzpicture}[
  start chain = A going right,
  node distance = 0pt, % Ensure no gap between bit nodes
  byte node/.style = {draw=lightgray, minimum height=20pt, text height=16pt, text depth=4pt, font=\ttfamily, anchor=south, align=center, inner sep=0pt, outer sep=0pt, on chain=A}
]
\newlength{\onebyte}
\setlength{\onebyte}{1.7ex}

\node[byte node,minimum width=4\onebyte,text width=4\onebyte,fill=white] (seg1) {@(\#)};
\node[byte node,minimum width=\onebyte,text width=\onebyte, fill=red!15,font=\scriptsize\ttfamily, text depth=0pt] (delim1a) {\rotatebox{90}{0xC2}};
\node[byte node,minimum width=\onebyte,text width=\onebyte, fill=red!15,font=\scriptsize\ttfamily, text depth=0pt] (delim1b) {\rotatebox{90}{0xA0} };
\node[byte node,minimum width=7\onebyte,text width=7\onebyte,fill=white] (seg2) {\emoji{robot}\#};
\node[byte node,minimum width=\onebyte,text width=\onebyte, fill=red!15,font=\scriptsize\ttfamily, text depth=0pt] (delim2a) {\rotatebox{90}{0xC2}};
\node[byte node,minimum width=\onebyte,text width=\onebyte, fill=red!15,font=\scriptsize\ttfamily, text depth=0pt] (delim2b) {\rotatebox{90}{0xA0} };
\node[byte node,minimum width=3\onebyte,text width=3\onebyte,fill=white] (seg3) {\emoji{white-check-mark}};
\node[byte node,minimum width=\onebyte,text width=\onebyte, fill=red!15,font=\scriptsize\ttfamily, text depth=0pt] (delim3a) {\rotatebox{90}{0xC2}};
\node[byte node,minimum width=\onebyte,text width=\onebyte, fill=red!15,font=\scriptsize\ttfamily, text depth=0pt] (delim3b) {\rotatebox{90}{0xA0} };
\node[byte node,minimum width=7\onebyte,text width=7\onebyte,fill=white] (seg4) {DwmWhat};
\node[byte node,minimum width=\onebyte,text width=\onebyte, fill=red!15,font=\scriptsize\ttfamily, text depth=0pt] (delim4a) {\rotatebox{90}{0xC2}};
\node[byte node,minimum width=\onebyte,text width=\onebyte, fill=red!15,font=\scriptsize\ttfamily, text depth=0pt] (delim4b) {\rotatebox{90}{0xA0} };
\node[byte node,minimum width=5\onebyte,text width=5\onebyte,fill=white] (seg5) {1.0.0};
\node[byte node,minimum width=\onebyte,text width=\onebyte, fill=red!15,font=\scriptsize\ttfamily, text depth=0pt] (delim5a) {\rotatebox{90}{0xC2}};
\node[byte node,minimum width=\onebyte,text width=\onebyte, fill=red!15,font=\scriptsize\ttfamily, text depth=0pt] (delim5b) {\rotatebox{90}{0xA0} };
\node[byte node,minimum width=18\onebyte,text width=18\onebyte,fill=white] (seg6) {Daniel McRobb 2025};
\node[byte node,minimum width=\onebyte,text width=\onebyte, fill=red!15,font=\scriptsize\ttfamily, text depth=0pt] (delim6a) {\rotatebox{90}{0xC2}};
\node[byte node,minimum width=\onebyte,text width=\onebyte, fill=red!15,font=\scriptsize\ttfamily, text depth=0pt] (delim6b) {\rotatebox{90}{0xA0} };
\node[byte node,minimum width=10\onebyte,text width=10\onebyte,fill=white] (seg7) {mcplex.net};
\node[byte node,minimum width=\onebyte,text width=\onebyte,fill=lightgray!50] (termnull) {$\emptyset$};

\node (nth0) [below=of seg1.south, yshift=-20pt,font=\scriptsize]
      {\texttt{nth(0)}};
\draw [decorate, decoration={brace, mirror, amplitude=10pt}]
  (seg2.south west) -- (seg2.south east)
  node [below=8pt, midway, font=\scriptsize, text=blue] {\texttt{type()}};
\node (nth1) [below=of seg2.south, yshift=-20pt,font=\scriptsize]
      {\texttt{nth(1)}};
\draw [decorate, decoration={brace, mirror, amplitude=10pt}]
  (seg3.south west) -- (seg3.south east)
  node [below=8pt, midway, font=\scriptsize, text=blue] {\texttt{status()}};
\node (nth2) [below=of seg3.south, yshift=-20pt,font=\scriptsize]
      {\texttt{nth(2)}};
\draw [decorate, decoration={brace, mirror, amplitude=10pt}]
  (seg4.south west) -- (seg4.south east)
  node [below=8pt, midway, font=\scriptsize, text=blue] {\texttt{name()}};
\node (nth3) [below=of seg4.south, yshift=-20pt,font=\scriptsize]
      {\texttt{nth(3)}};
\draw [decorate, decoration={brace, mirror, amplitude=10pt}]
  (seg5.south west) -- (seg5.south east)
  node [below=8pt, midway, font=\scriptsize, text=blue] {\texttt{version()}};
\node (nth4) [below=of seg5.south, yshift=-20pt,font=\scriptsize]
      {\texttt{nth(4)}};
\draw [decorate, decoration={brace, mirror, amplitude=10pt}]
  (seg6.south west) -- (seg6.south east)
  node [below=8pt, midway, font=\scriptsize, text=blue] {\texttt{copyright()}};
\node (nth6) [below=of seg6.south, yshift=-20pt,font=\scriptsize]
      {\texttt{nth(5)}};
\draw [decorate, decoration={brace, mirror, amplitude=10pt}]
  (seg7.south west) -- (seg7.south east)
  node [below=8pt, midway, font=\scriptsize, text=blue] {other()};
\node (nth8) [below=of seg7.south, yshift=-20pt,font=\scriptsize]
      {\texttt{nth(6)}};
      
\draw [decorate, decoration={brace, amplitude=10pt}]
  (seg2.north west) -- (seg7.north east)
  node [above=8pt, midway, font=\scriptsize, text=blue] {\texttt{data\_view()}};

\draw [decorate, decoration={brace, amplitude=40pt}]
  (seg1.north west) -- (seg7.north east)
  node [above=40pt, midway, font=\scriptsize] {\texttt{view()}};

\end{tikzpicture}
\caption{\texttt{Dwm::What::Info} layout}
\label{figure:InfoEx01}
\end{figure}
\end{center}
\end{minipage}

The 67 bytes:
\begin{tabular}{ll}
  @(\#)            &  4 bytes \\
  delim            &  2 bytes \\
  type             &  7 bytes \\
  delim            &  2 bytes \\
  status           &  3 bytes \\
  delim            &  2 bytes \\
  name             &  7 bytes \\
  delim            &  2 bytes \\
  version          &  5 bytes \\
  delim            &  2 bytes \\
  copyright        & 18 bytes \\
  delim            &  2 bytes \\
  other            & 10 bytes \\
  null             &  1 byte \\ \hline
  \textbf{Total} & \textbf{67 bytes}
\end{tabular}

And {\em dwmwhat} would print this when scanning a compiled program
containing such an instance of \texttt{Dwm::What::Info}:

\texttt{
\emoji{robot}\# \emoji{white-check-mark} DwmWhat 1.0.0 Daniel McRobb 2025 mcplex.net
}
